\chapter{代数结构}\label{ch:b1:structures}

同样的运算规则反复出现：整数、实数、
复数、\omterm{def:b1:arith:congruence}{同余}类，以及即将学到的多项式
（\cref{ch:b1:poly}）、向量与矩阵
（\cref{ch:b1:vspaces,ch:b1:matrices}）。代数抽取其中的
共同规律，并给它们命名：\emph{群}、\emph{环}、\emph{域}。
在结构层面证明一次某条事实，
便对所有实例同时成立。

\section{运算}

\begin{definition}\label{def:b1:structures:law}
一个\emph{运算}\index{运算}，是\omterm{def:b1:logic:sets}{集合} $E$ 上的
\omterm{def:b1:logic:map}{映射} $E \times E \to E$，记作 $(x, y) \mapsto x * y$。若运算是
\emph{结合的}，即恒有 $(x*y)*z = x*(y*z)$，
运算是\emph{交换的}，即恒有 $x * y = y * x$。元素 $e$ 称为\emph{单位元}，若
$e * x = x * e = x$ 对所有 $x$ 成立；此时 $x'$ 称为 $x$ 的\emph{逆元}，若
$x * x' = x' * x = e$。
\end{definition}

\begin{proposition}[唯一性]\label{prop:b1:structures:unique}
一个运算至多有一个单位元；对于带单位元的结合运算，
每个元素至多有一个逆元。
\end{proposition}

\begin{proof}
若 $e$ 与 $e'$ 都是单位元，则 $e = e * e' = e'$。若 $x'$ 与 $x''$
都是 $x$ 的逆元，则 $x' = x' * e = x' * (x * x'') = (x' * x) * x'' = e * x''
= x''$。
\end{proof}

\section{群}

\begin{definition}[群]\label{def:b1:structures:group}
一个\emph{群}\index{群} $(G, *)$ 是指一个\omterm{def:b1:logic:sets}{集合}，配以满足结合律的
运算，且存在单位元，每个元素都有逆元。
若该运算还满足
交换律，则称该群为\emph{阿贝尔群}\index{阿贝尔群}。
\end{definition}

\begin{example}\label{ex:b1:structures:groups}
$(\Z, +)$, $(\Q, +)$, $(\R, +)$, $(\C, +)$; $(\Q^*, \times)$,
$(\R^*, \times)$，$(\C^*, \times)$，$(\mathbb{U}_n, \times)$（\omterm{def:b1:complex:unity}{单位根}，
见\cref{def:b1:complex:unity}）；记 $\mathfrak{S}(E)$ 为
\omterm{def:b1:logic:sets}{集合} $E$ 到自身的全体\omterm{def:b1:logic:inj}{双射}，在复合运算下构成的群，
称为\emph{对称群}\index{对称群}（$E$ 上的），
当 $\abs E \geq 3$ 时它便不是\omterm{def:b1:structures:group}{阿贝尔群}。不构成群的例子有 $(\N, +)$（无逆元），
$(\Z, \times)$（只有 $\pm 1$ 可逆）。
\end{example}

\begin{proposition}[运算规则]\label{prop:b1:structures:rules}
设 $G$ 为群（以乘法记号书写，单位元为 $e$）：
\begin{enumerate}
  \item 消去律：$ax = ay \implies x = y$ 且 $xa = ya \implies
        x = y$；
  \item $(ab)^{-1} = b^{-1} a^{-1}$ 且 $(a^{-1})^{-1} = a$；
  \item 对 $a, b \in G$，方程 $ax = b$ 与 $xa = b$ 各有
        唯一解（分别为 $x = a^{-1}b$，$x = b a^{-1}$）。
\end{enumerate}
\end{proposition}

\begin{proof}
（1）在适当一侧乘以 $a^{-1}$，并利用
结合律。（2）$(b^{-1}a^{-1})(ab) = b^{-1}(a^{-1}a)b = b^{-1}b =
e$，对称地可得；由逆元的唯一性得出结论；第二条
由\cref{prop:b1:structures:unique}对 $a^{-1}$ 应用即得。（3）
代入后利用（1）得唯一性。
\end{proof}

\begin{example}[矩形的对称性]\label{ex:b1:structures:klein}
非正方形的矩形恰有四个到自身的等距变换，它们是
恒等变换 $e$、水平轴反射 $h$、
竖直轴反射 $v$，以及绕中心的半转 $r$。
复合运算使这个四元\omterm{def:b1:logic:sets}{集合}构成群：每个
元素都是自己的逆元（$h^2 = v^2 = r^2 = e$），并且
任意两个不同的非单位元之积恰为第三个非单位元，
（$hv = vh = r$：沿两轴反射即为半转）。完整的
乘法表对称，故该群为\omterm{def:b1:structures:group}{阿贝尔群}——然而它
\emph{不}同于 \cref{ex:b1:structures:order}中的旋转群 $\mathbb U_4$：
那里 $\iu$ 的阶为 $4$，而
这里则每个元素的阶都 $\leq 2$。元素个数相同的两个群
因此可以具有本质不同的乘法结构，
——下图并排展示两张表。这个四元群
又可写作 $\{\pm1\} \times \{\pm1\}$，并
由\cref{exo:b1:structures:7}可知，所有元素的平方都是单位元的群
必定像它一样是\omterm{def:b1:structures:group}{阿贝尔群}。
\end{example}

\begin{omfigure}
\begin{tikzpicture}[scale=0.62]
  \foreach \i/\l in {0/$e$, 1/$\iu$, 2/$-1$, 3/$-\iu$} {
    \node[font=\small] at ({\i + 0.5}, 0.5) {\l};
    \node[font=\small] at (-0.5, {-\i - 0.5}) {\l};
  }
  \fill[omDef!25] (0,-1) rectangle ++(1,1);
  \fill[omDef!25] (3,-2) rectangle ++(1,1);
  \fill[omDef!25] (2,-3) rectangle ++(1,1);
  \fill[omDef!25] (1,-4) rectangle ++(1,1);
  \draw[gray] (0,-4) grid (4,0);
  \foreach \i/\j/\l in {0/0/$e$, 0/1/$\iu$, 0/2/$-1$, 0/3/$-\iu$,
    1/0/$\iu$, 1/1/$-1$, 1/2/$-\iu$, 1/3/$e$,
    2/0/$-1$, 2/1/$-\iu$, 2/2/$e$, 2/3/$\iu$,
    3/0/$-\iu$, 3/1/$e$, 3/2/$\iu$, 3/3/$-1$}
    \node[font=\scriptsize] at ({\j+0.5}, {-\i-0.5}) {\l};
  \begin{scope}[xshift=7.5cm]
    \foreach \i/\l in {0/$e$, 1/$h$, 2/$v$, 3/$r$} {
      \node[font=\small] at ({\i + 0.5}, 0.5) {\l};
      \node[font=\small] at (-0.5, {-\i - 0.5}) {\l};
    }
    \fill[omThm!25] (0,-1) rectangle ++(1,1);
    \fill[omThm!25] (1,-2) rectangle ++(1,1);
    \fill[omThm!25] (2,-3) rectangle ++(1,1);
    \fill[omThm!25] (3,-4) rectangle ++(1,1);
    \draw[gray] (0,-4) grid (4,0);
    \foreach \i/\j/\l in {0/0/$e$, 0/1/$h$, 0/2/$v$, 0/3/$r$,
      1/0/$h$, 1/1/$e$, 1/2/$r$, 1/3/$v$,
      2/0/$v$, 2/1/$r$, 2/2/$e$, 2/3/$h$,
      3/0/$r$, 3/1/$v$, 3/2/$h$, 3/3/$e$}
      \node[font=\scriptsize] at ({\j+0.5}, {-\i-0.5}) {\l};
  \end{scope}
\end{tikzpicture}

{\small 两个四元群：$\mathbb U_4 = \{e, \iu,
-1, -\iu\}$（左）与矩形群（右），其中
单位元的位置涂有阴影。左图中，单位元呈蜿蜒状
（一个 $4$ 阶元即可生成整个群）；右图中，它
填满对角线（每个元素的平方都是 $e$）。无论怎样重新标记，
都无法把一张表变成另一张：两个群
不同构。}
\end{omfigure}

\begin{definition}[子群]\label{def:b1:structures:subgroup}
$H$ 为群 $G$ 的子集，若它是\emph{子群}\index{子群}
（记作 $H \leq G$），它须含有 $e$，且对
运算与求逆封闭。此时 $H$ 本身也是群。

\emph{判别准则：}非空子集 $H \subseteq G$ 是子群，当且
仅当下式成立：
\[
\forall x, y \in H, \quad x y^{-1} \in H .
\]
\end{definition}

\begin{proof}[准则的证明]
\omterm{def:b1:structures:subgroup}{子群}显然满足该条件。反之，设 $H \neq \emptyset$
且满足该条件，取 $x_0 \in H$。于是 $e = x_0 x_0^{-1} \in H$；对任意
$y \in H$，有 $y^{-1} = e\,y^{-1} \in H$；对任意 $x, y \in H$，$xy =
x (y^{-1})^{-1} \in H$。
\end{proof}

\begin{example}\label{ex:b1:structures:subgroups}
$\mathbb{U}_n \leq (\C^*, \times)$：非空，且对 $z, w \in
\mathbb{U}_n$，$(zw^{-1})^n = z^n (w^n)^{-1} = 1$。$(\Z, +)$ 的\omterm{def:b1:structures:subgroup}{子群}
恰为 $n\Z$（证明见
\cref{thm:b1:arith:gcd}）。\omterm{def:b1:structures:subgroup}{子群}的交总是
\omterm{def:b1:structures:subgroup}{子群}，但并集几乎从不是
（见\cref{exo:b1:structures:6}）。
\end{example}

\begin{remark}[结构的常见误区]\label{rem:b1:structures:pitfalls}
\begin{enumerate}
  \item \emph{仅对运算封闭是不够的。}$\N$ 在
        $\Z$ 中对加法封闭，含有 $0$，却
        不是\omterm{def:b1:structures:subgroup}{子群}：缺少逆元。判别准则
        $xy^{-1} \in H$ 一次检验全部条件——但只有
        在检验 $H \neq \emptyset$ 之后才行。
  \item \emph{非阿贝尔情形的思维定式。}在一般群中，$(ab)^2 =
        abab$，它等于 $a^2b^2$，而这仅当 $a$ 与 $b$ 可交换时成立；
        同理，$(ab)^{-1} = b^{-1}a^{-1}$，次序颠倒。
        由中学代数移植来的恒等式，都须
        从公理重新推导，或标明仅适用于交换情形。
  \item \emph{核与像的区别。}$\ker f$ 位于
        \emph{源群}中，$\operatorname{im} f$ 位于目标群中；
        “$f$ 为\omterm{def:b1:logic:inj}{单射}当且仅当 $\ker f$ 平凡”
        （见\cref{prop:b1:structures:kernel}）对像没有类似的结论，
        （$\operatorname{im} f = G'$ 即\omterm{def:b1:logic:inj}{满射}性）。
  \item \emph{环关于 $\times$ 不是群。}在环中，多数
        元素未必可逆，而用 $a$ 消去时
        要求 $a$ 为可逆元，或环为\omterm{def:b1:structures:field}{整
        环}。例如在 $\Z/12\Z$ 中，$\overline3\,\overline2 =
        \overline3\,\overline6$，但 $\overline2 \neq \overline6$
        （见\cref{ex:b1:structures:zncomputation}）。
\end{enumerate}
\end{remark}

\begin{definition}[群同态]\label{def:b1:structures:morphism}
设 $(G, *)$ 与 $(G', \star)$ 为群，\omterm{def:b1:logic:map}{映射} $f \colon G \to G'$
为\emph{同态}\index{同态!群的}，当
\[
\forall x, y \in G, \qquad f(x * y) = f(x) \star f(y).
\]
此时 $f(e_G) = e_{G'}$ 且 $f(x^{-1}) = f(x)^{-1}$。
$f$ 的\emph{核}\index{核}与\emph{像}为
\[
\ker f = f^{-1}(\{e_{G'}\}) \leq G,
\qquad
\operatorname{im} f = f(G) \leq G' .
\]
\omterm{def:b1:logic:inj}{双射}同态称为\emph{同构}；此时其逆\omterm{def:b1:logic:map}{映射}
自动是同态。
\end{definition}

\begin{proof}[论断的证明]
$f(e) = f(e * e) = f(e)\star f(e)$，消去 $f(e)$ 得
$e_{G'} = f(e)$。于是 $f(x)\star f(x^{-1}) = f(x x^{-1}) = e_{G'}$
由此 $f(x^{-1})$ 即为其逆元。核：$e \in \ker f$；若
$x, y \in \ker f$，$f(xy^{-1}) = f(x)f(y)^{-1} = e$；由
判别准则即得。像：同样由判别准则，并利用 $f(x)f(y)^{-1} = f(xy^{-1})$。
同构的逆：对 $u, v \in G'$，记 $u = f(x)$，$v =
f(y)$，则 $f^{-1}(u \star v) = f^{-1}(f(xy)) = xy =
f^{-1}(u) f^{-1}(v)$。
\end{proof}

\begin{proposition}[借助核判定单射性]\label{prop:b1:structures:kernel}
群\omterm{def:b1:structures:morphism}{同态} $f$ 是\omterm{def:b1:logic:inj}{单射}，当且仅当 $\ker f =
\{e\}$。
\end{proposition}

\begin{proof}
若 $f$ 是\omterm{def:b1:logic:inj}{单射}，则 $\ker f$ 中只能含有
$e_{G'}$ 的唯一\omterm{def:b1:logic:map}{原像}，即 $e$。反之，若 $\ker f = \{e\}$ 且 $f(x) =
f(y)$，则 $f(xy^{-1}) = f(x) f(y)^{-1} = e_{G'}$，从而 $xy^{-1} = e$，
即 $x = y$。
\end{proof}

\begin{example}\label{ex:b1:structures:expmorphism}
$\exp \colon (\R, +) \to (\R_+^*, \times)$ 是\omterm{def:b1:structures:morphism}{同态}
（$\eu^{x+y} = \eu^x \eu^y$），且是\omterm{def:b1:logic:inj}{双射}
（见\cref{prop:b1:functions:expln}）：加法结构与乘法结构
同构——这正是对数在历史上
存在的理由。另一个\omterm{def:b1:structures:morphism}{同态}：$\theta \mapsto \eu^{\iu\theta}$ 从
$(\R, +)$ \omterm{def:b1:logic:inj}{满射}到单位圆 $(\mathbb{U}, \times)$，核为
$2\pi\Z$。
\end{example}

\begin{example}[符号同态]\label{ex:b1:structures:signmorphism}
\omterm{def:b1:logic:map}{映射} $s \colon (\R^*, \times) \to (\{\pm1\}, \times)$ 把
$x$ 映为其符号，是一个\omterm{def:b1:structures:morphism}{同态}：乘积的符号
等于各因子符号之积。其核为 $\intoo0{+\infty}$（一个\omterm{def:b1:structures:subgroup}{子群}，
正如\cref{def:b1:structures:morphism}所许诺的），其像
恰为全部 $\{\pm1\}$：\omterm{def:b1:logic:inj}{满射}，而且极度非\omterm{def:b1:logic:inj}{单射}。
两条一般性的道理，在此以小见大。其一，\omterm{def:b1:structures:morphism}{同态}可能抹去大量信息：
$s$ 对 $x$ 只保留一位信息，这正是它的长处——符号论证恰好是
那些经由 $s$ 分解的计算。其二，
到 $\{\pm1\}$ 的\omterm{def:b1:structures:morphism}{同态}是最简单的“不变量”：
即本章周末问题中构造的\omterm{def:b1:counting:objects}{置换}符号，
在\omterm{ex:b1:structures:groups}{对称群} $\mathfrak S_n$ 上是同一现象，
而其所依赖的奇偶性论证，
都经由这样的两值\omterm{def:b1:structures:morphism}{同态}导出。
\end{example}

\begin{definition}[幂，元素的阶]\label{def:b1:structures:order}
在群 $G$ 中（采用乘法记号），令 $x^0 = e$，$x^{k+1} =
x^k x$，$x^{-k} = (x^k)^{-1}$，其中 $k \in \N$；于是 $x^{k+l} = x^k
x^l$（对一切 $k, l \in \Z$ 成立），因此 $k \mapsto x^k$ 是 $(\Z,
+) \to G$ 的\omterm{def:b1:structures:morphism}{同态}，其像 $\langle x \rangle = \{x^k : k \in \Z\}$ 是
一个\omterm{def:b1:structures:subgroup}{子群}，即由 $x$ \emph{生成}的\omterm{def:b1:structures:subgroup}{子群}。
$x$ 的\emph{阶}\index{阶}是使 $m \geq
1$ 且 $x^m = e$ 的最小者（若存在；此时 $\langle x\rangle = \{e, x,
\dots, x^{m-1}\}$ 恰有 $m$ 个元素，且 $x^k = e \iff m
\mid k$），否则阶为 $\infty$。
\end{definition}

\begin{example}\label{ex:b1:structures:order}
在 $(\C^*, \times)$ 中：$\iu$ 的阶为 $4$，且 $\langle \iu \rangle
= \{1, \iu, -1, -\iu\} = \mathbb{U}_4$；更一般地，$\omega =
\eu^{2\iu\pi/n}$ 的阶为 $n$，且 $\langle\omega\rangle =
\mathbb{U}_n$。在 $(\Z, +)$ 中，每个 $x \neq 0$ 的阶均为无穷。
定义中各断言成立的理由：若 $x$ 的阶为 $m$，
对任一 $k$ 以 $m$ 作带余除法（$k = mq + r$，$0 \leq r < m$，
\cref{thm:b1:arith:division}）：$x^k = (x^m)^q x^r = x^r$，于是
幂以 $m$ 为周期循环，所列元素由 $m$ 的极小性两两不同，
而 $x^k = e$ 迫使 $r = 0$。\omterm{def:b1:counting:objects}{置换}的阶
在下面的周末问题中计算。
\end{example}

\begin{example}[$\mathbb U_{12}$ 中的阶]\label{ex:b1:structures:orderU12}
$\omega^k$ 在 $\mathbb U_n$ 中的阶是多少？其中 $\omega =
\eu^{2\iu\pi/n}$？即有 $(\omega^k)^m = 1$ 当且仅当 $n \mid km$，且
记 $d = \gcd(n, k)$，$n = dn'$，$k = dk'$，且 $\gcd(n', k')
= 1$：$n \mid km \iff n' \mid k'm \iff n' \mid m$（高斯
引理，\cref{thm:b1:arith:gauss}）。满足条件的最小 $m \geq 1$ 为
$n' = \frac{n}{\gcd(n,k)}$。例如在 $\mathbb U_{12}$ 中，
$\omega^8$ 的阶为 $\frac{12}{\gcd(12,8)} = 3$（事实上
$\omega^8 = \eu^{4\iu\pi/3} \in \mathbb U_3$），而 $\omega^5$
的阶为 $12$：它生成整个群，尽管它不是
“标准”的生成元。计数生成元（即 $k$
满足 $\gcd(k, n) = 1$ 的那些）——恰好复现
\cref{ex:b1:counting:coprime}中的\omterm{cor:b1:arith:bezout}{互质}计数：群论与计数在此相遇。
\end{example}

\section{环与域}

\begin{definition}[环]\label{def:b1:structures:ring}
\emph{环}\index{环} $(A, +, \times)$ 是带两种运算的\omterm{def:b1:logic:sets}{集合}，使得
$(A, +)$ 是\omterm{def:b1:structures:group}{阿贝尔群}（单位元为 $0$）；$\times$ 满足结合律，
且有单位元 $1$；并且 $\times$ 对 $+$ 的分配律在两侧成立。
当 $\times$ 满足交换律时，该环称为\emph{交换的}。一个元素 $a$ 称为
\emph{可逆的}（亦称\emph{单位}），当 $ab = ba = 1$
对某个 $b$ 成立时；全体单位构成群 $(A^\times, \times)$。
\end{definition}

\begin{proof}[证明单位构成群]
稳定性：若 $a, a'$ 为单位，其逆分别为 $b, b'$，则
\[
(aa')(b'b) = a(a'b')b = a\,1\,b = ab = 1,
\qquad (b'b)(aa') = 1
\]
同理可知 $aa'$ 为单位。元素 $1$ 是单位（其
自身的逆），结合律继承自 $A$，而
$b$ 为单位 $a$ 的逆，其自身亦为单位（其逆为 $a$）。
于是 $(A^\times, \times)$ 满足群的全部公理。本书中凡非由\omterm{def:b1:counting:objects}{置换}构成的群，
都以此方式产生：
$\Q^* = \Q^\times$、$\R^*$、$\C^*$，以及下文的 $\Z/n\Z$ 的单位群，
以及以后的可逆矩阵（见\cref{ch:b1:matrices}）。
\end{proof}

\begin{example}\label{ex:b1:structures:rings}
$\Z, \Q, \R, \C$ 都是交换环；$\Z^\times = \{1, -1\}$，
$\Q^\times = \Q^*$。以后还将遇到：多项式环 $K[X]$
（见\cref{ch:b1:poly}）、矩阵环（不可交换，
见\cref{ch:b1:matrices}），以及下文的 $\Z/n\Z$。在任一环中，$0 \times
a = 0$（由分配律：$0a = (0+0)a = 0a + 0a$），且
$(-1)a = -a$.
\end{example}

\begin{example}[幂等元：新环中的新现象]\label{ex:b1:structures:idempotents}
在 $\Z$ 中，方程 $x^2 = x$，即 $x(x - 1) = 0$，仅有
解 $0$ 与 $1$。在 $\Z/6\Z$ 中逐一检验所有剩余类：
$\overline0^2 = \overline0$, $\overline1^2 = \overline1$,
$\overline3^2 = \overline9 = \overline3$ 与 $\overline4^2 =
\overline{16} = \overline4$——共\emph{四}个幂等元。其中两个
奇异者来自零因子：$\overline3\,(\overline3 -
\overline1) = \overline3 \times \overline2 = \overline6 =
\overline0$，而两个因子都不为零。此类计算有助于
校准直觉：关于方程的熟悉事实在
\omterm{def:b1:structures:field}{整环}与域中仍然成立，但一般的环可以、并且确实会
表现出不同的性质——另见
\cref{exo:b1:structures:10}中的布尔环，其中\emph{每个}元素都
是幂等的。
\end{example}

\begin{proposition}[交换环中的二项式定理]\label{prop:b1:structures:binomial}
若 $a, b$ 是交换环中的元素（更一般地，若
$ab = ba$），则对 $n \in \N$：
\[
(a+b)^n = \sum_{k=0}^n \binom nk a^k b^{n-k},
\qquad
a^n - b^n = (a - b) \sum_{k=0}^{n-1} a^k b^{\,n-1-k} .
\]
\end{proposition}

\begin{proof}
\cref{thm:b1:counting:binomial}的证明及几何恒等式的
证明只用到结合律、两元素的交换律，
以及分配律——它们可原样套用。
\end{proof}

\begin{example}[陌生环中的二项式定理]\label{ex:b1:structures:binomialring}
这种一般性带来两个快捷的推论。在 $\Z/p\Z$ 中（$p$ 为\omterm{def:b1:arith:prime}{质数}），
中间的\omterm{def:b1:counting:objects}{二项式系数}全部为零
（\cref{thm:b1:arith:fermat}的第一步），于是定理化为
所谓的\emph{新生之梦}
\[
(a + b)^p = a^p + b^p \qquad \text{在 } \Z/p\Z ,
\]
在该环中它是一个真正的恒等式，尽管在 $\R$ 上看起来十分荒谬。
又若某交换环中含有元素 $\varepsilon$，
满足 $\varepsilon^2 = 0$，则定理在此截断：
$(a + \varepsilon)^n = a^n + n\,a^{n-1}\varepsilon$，所有更高次项
都含有因子 $\varepsilon^2 = 0$。
$n\,a^{n-1}$ 为 $\varepsilon$ 的系数，恰是 $x^n$ 的导数——
这绝非偶然，它是一个初步的暗示：导数
既属于分析，也同样属于代数（比较
\cref{ch:b1:poly}中的形式导数）。
\end{example}

\begin{definition}[整环、域]\label{def:b1:structures:field}
称交换环 $A \neq \{0\}$ 为\emph{整环}\index{整环}，当它
无零因子时，即 $ab = 0
\implies a = 0$ 或 $b = 0$。若每个非零元素都可逆，则称其为\emph{域}\index{域}。
每个域都是整环
（由 $ab = 0$ 及 $a \neq 0$ 得 $b = a^{-1}ab = 0$）。
\end{definition}

\begin{example}
$\Q$、$\R$、$\C$ 都是域；$\Z$ 是\omterm{def:b1:structures:field}{整环}，但不是
域。在\omterm{def:b1:structures:field}{整环}中，关于 $\times$ 的消去律成立：$ab =
ac$ 且 $a \neq 0$ 蕴含 $b = c$。
\end{example}

\section{环 $\Z/n\Z$}

\begin{definition}\label{def:b1:structures:zn}
取定 $n \in \N^*$。模 $n$ 的\omterm{def:b1:arith:congruence}{同余}类
（见\cref{ex:b1:logic:congruence}）构成一个
\omterm{def:b1:logic:sets}{集合} $\Z/n\Z$\index{Z/nZ@$\Z/n\Z$}，含 $n$ 个元素，记为 $\overline 0,
\overline 1, \dots, \overline{n-1}$。运算
\[
\overline a + \overline b = \overline{a + b},
\qquad
\overline a \times \overline b = \overline{ab}
\]
是良定义的——结果的\omterm{def:b1:arith:congruence}{同余}类不依赖于
代表元的选取，正是因为\omterm{def:b1:arith:congruence}{同余}关于 $+$
与 $\times$ 相容（见\cref{def:b1:arith:congruence}）——并使 $\Z/n\Z$ 成为
交换环。
\end{definition}

\begin{theorem}[$\Z/n\Z$ 的可逆元；域 $\Z/p\Z$]\label{thm:b1:structures:zp}
\begin{enumerate}
  \item $\overline a$ 在 $\Z/n\Z$ 中可逆当且仅当
        $\gcd(a, n) = 1$。
  \item $\Z/n\Z$ 是域当且仅当 $n$ 为\omterm{def:b1:arith:prime}{质数}。
\end{enumerate}
\end{theorem}

\begin{proof}
(1) 即用\omterm{def:b1:arith:congruence}{同余}类的语言重述 \cref{prop:b1:arith:invmod}。

(2) 若 $n = p$ 为\omterm{def:b1:arith:prime}{质数}，则每个 $\overline a \neq \overline 0$ 满足 $p
\nmid a$，故 $\gcd(a, p) = 1$，由 (1) 知其可逆——从而是域。若 $n =
ab$，其中 $1 < a, b < n$，则 $\overline a\, \overline b = \overline
n = \overline 0$，且 $\overline a, \overline b \neq \overline 0$：
它们是零因子，故甚至不是\omterm{def:b1:structures:field}{整环}；而 $n = 1$ 给出
零环，应予排除。
\end{proof}

\begin{example}[$1$ 有几个平方根？]\label{ex:b1:structures:sqrtone}
求 $x^2 = \overline 1$ 在 $\Z/8\Z$ 与 $\Z/7\Z$ 中的解。逐一检验模 $8$ 的
八个\omterm{def:b1:arith:congruence}{同余}类：$1^2 = 1$，$3^2 = 9 \equiv 1$，$5^2 = 25
\equiv 1$，$7^2 = 49 \equiv 1$——\emph{四个}解
$\{\overline1, \overline3, \overline5, \overline7\}$，尽管
多项式 $X^2 - 1$ 的次数为 $2$。在域 $\Z/7\Z$ 中，
相反，$x^2 = \overline1$ 等价于 $(x - \overline1)(x +
\overline1) = \overline0$，而域中没有零因子：$x =
\pm\overline1$，只有两个解。模 $8$ 时的失效
是可以追溯的：$(3-1)(3+1) = 2 \times 4 = 8 \equiv 0$，而两个因子
均不为零。由此得教训：熟悉的规则“$d$ 次方程
至多有 $d$ 个根”是关于
\emph{\omterm{def:b1:structures:field}{整环}}（\cref{cor:b1:poly:nroots}在域上证明了它）；
而在含零因子的环中它悄然失效——这恰好说明
为何威尔逊定理的配对证明
（\cref{exo:b1:arith:11}）需要 $p$ 为\omterm{def:b1:arith:prime}{质数}。
\end{example}

\begin{example}[在 $\Z/n\Z$ 中计算]\label{ex:b1:structures:zncomputation}
在 $\Z/12\Z$ 中：可逆元为 $\overline 1, \overline 5, \overline 7,
\overline{11}$（即与 $12$ \omterm{cor:b1:arith:bezout}{互质}的\omterm{def:b1:arith:congruence}{同余}类），且每个元素都是自己的
逆元（$5^2 = 25 \equiv 1$，$7^2 = 49 \equiv 1$，$11^2 = 121 \equiv
1$）。方程 $\overline 3\, x = \overline 6$ 有\emph{三个}
解（$x \in \{\overline 2, \overline 6, \overline{10}\}$）：
若无可逆性，便不能消去。相反，在 $\Z/11\Z$ 中，
每个形如 $\overline a x = \overline b$（$\overline a \neq
\overline 0$）的方程恰有一个解。
\end{example}

\begin{example}[群公理作为求解的保证]\label{ex:b1:structures:solving}
在群 $\bigl((\Z/7\Z)^*, \times\bigr)$ 中，求解 $\overline
3\,x = \overline 5$。由\cref{prop:b1:structures:rules}~(3)，
解存在、唯一，且等于 $\overline3^{-1}\,
\overline5$；由于 $\overline3 \times \overline5 = \overline{15}
= \overline1$，故 $\overline 3$ 的逆元为 $\overline 5$，因此
\[
x = \overline5 \times \overline5 = \overline{25} = \overline4,
\qquad\text{ 验算： } \overline3 \times \overline4 =
\overline{12} = \overline5 .
\]
重点不在于答案，而在于保证：在群中，
每个这样的方程都唯一可解，\emph{先于}任何
计算即已确定，故求解过程绝不会遇到“无
解”或“多解”。对比上文中 $\overline3\,x = \overline6$
在 $\Z/12\Z$ 中的情形，那里保证失效——知道
身处何种结构，便是知道可以视为理所当然的是什么。
\end{example}

\begin{example}[直积]\label{ex:b1:structures:product}
设 $G$ 与 $H$ 为群，则笛卡儿积 $G \times H$ 配以
逐分量的运算 $(g, h)(g', h') = (gg', hh')$ 构成一个群：
公理逐个分量验证，单位元为
$(e_G, e_H)$，逆元为 $(g^{-1}, h^{-1})$。阶按
\omterm{def:b1:arith:lcm}{最小公倍数}合成：$(g, h)^m = (g^m, h^m)$ 为单位元当且仅当
$g$ 的阶与 $h$ 的阶都\omterm{def:b1:arith:divides}{整除} $m$。于是在 $\Z/2\Z
\times \Z/2\Z$（加法）中每个非零元的阶都是 $2$
——这恰是
\cref{ex:b1:structures:klein} 中矩形群在坐标下的表示——而 $\Z/4\Z$
有一个阶为 $4$ 的元素：这是另一种无需计算的证明，
说明两个 $4$ 元群不同构（同构
保持阶）。直积是由已知群构造新群最简便的方法，
而平面 $\R^2 = \R \times
\R$（见\cref{ch:b1:vspaces}）是该构造
最重要的实例。
\end{example}

\begin{remark}[从结构看费马定理]
在域 $\Z/p\Z$ 中，非零剩余类构成一个乘法群，
其元素个数为 $p - 1$，而费马小定理
（\cref{thm:b1:arith:fermat}）断言：该群的每个元素 $x$
满足 $x^{p-1} = \overline 1$。这是有限群上一般事实的特例，
即拉格朗日定理，将在本科
二年级予以证明。威尔逊定理的配对证明
（\cref{exo:b1:arith:11}）已具有这种群论色彩。
\end{remark}

\begin{remark}[插曲：抽象带来了什么]\label{rem:b1:structures:interlude}
不妨问一问，对抽象的运算而非对具体的数证明
\cref{prop:b1:structures:unique}，究竟有何收获？
答案在于杠杆作用。这一两行的论证如今
可同时涵盖：复合运算下函数的逆，
（\cref{thm:b1:logic:inverse}，其唯一性证明
它一字不差地重复）、模 $n$ 的逆、
（\cref{prop:b1:arith:invmod}）、非零实数的逆、
任意环中单位的逆，以及——未经细看——在
\cref{ch:b1:matrices}中的可逆矩阵的逆，其中 $A^{-1}$ 的唯一性
无需一行证明。同样的节约也适用于
\cref{prop:b1:structures:kernel}（一条\omterm{def:b1:logic:inj}{单射}判别准则，
在 \cref{ch:b1:linmaps} 中被线性\omterm{def:b1:logic:map}{映射}复用）和
\omterm{def:b1:structures:subgroup}{子群}判别法。这里的抽象并非为泛化而泛化：
它的实质是拒绝把同一个引理证明五遍，
在五个名字下各证一遍。代价是弄清每个陈述
究竟用了哪些公理——这正是本章习题
所训练的。
\end{remark}

\begin{remark}[本章的应用]\label{rem:b1:structures:whereused}
本章的词汇正是本册其余部分的语法。环与域
组织了 \cref{ch:b1:poly}（$K[X]$ 是
模仿 $\Z$ 的环）和 \cref{ch:b1:fractions}（$K(X)$ 是其
分式域）；向量空间（\cref{ch:b1:vspaces}）是
带有域作用其上的\omterm{def:b1:structures:group}{阿贝尔群}；矩阵
（\cref{ch:b1:matrices}）构成本册第一个真正非交换的环，
而其可逆元构成一个群，
其研究本身就是线性代数。\omterm{def:b1:structures:morphism}{同态}与核将以
线性\omterm{def:b1:logic:map}{映射}与零空间的面貌在 \cref{ch:b1:linmaps} 中重现——
\cref{prop:b1:structures:kernel} \emph{正是}该章的\omterm{def:b1:logic:inj}{单射}判别准则，
在此已一劳永逸地证明。
下面周末问题的主角\omterm{ex:b1:structures:groups}{对称群}，提供了
行列式在\cref{ch:b1:det}中赖以建立的\omterm{def:b1:counting:objects}{置换}符号。
\end{remark}

\section{练习}

\begin{exercise}[$\star$]\label{exo:b1:structures:1}
在 $E = \R \setminus \{1\}$ 上定义 $x * y = x + y - xy$，证明
$(E, *)$ 是\omterm{def:b1:structures:group}{阿贝尔群}。\emph{（求出单位元及
$x$ 的逆元；验证封闭性：为何 $x * y \neq 1$？）}
\end{exercise}

\begin{exercise}[$\star$]\label{exo:b1:structures:2}
下列哪些是群？
\begin{enumerate}
  \item $(\intoo{0}{+\infty}, \times)$;
  \item $(\{-1, 0, 1\}, +)$;
  \item $(\Q^*, \times)$;
  \item 奇整数全体构成的\omterm{def:b1:logic:sets}{集合}，运算为加法。
\end{enumerate}
\end{exercise}

\begin{exercise}[$\star$]\label{exo:b1:structures:3}
写出\omterm{ex:b1:structures:groups}{对称群} $\mathfrak{S}_3$ 的运算表
（即 $\{1,2,3\}$ 上的六个\omterm{def:b1:logic:inj}{双射}：恒等\omterm{def:b1:logic:map}{映射}、三个对换、两个
$3$-轮换），并举出两个不可交换的元素。
\end{exercise}

\begin{exercise}[$\star$]\label{exo:b1:structures:4}
证明 $H = \{z \in \C^* : \abs z = 1\}$ 是 $(\C^*,
\times)$ 的\omterm{def:b1:structures:subgroup}{子群}，$\R_+^*$ 也是另一个\omterm{def:b1:structures:subgroup}{子群}；$H \cup \R_+^*$ 是否
仍为\omterm{def:b1:structures:subgroup}{子群}？
\end{exercise}

\begin{exercise}[$\star\star$]\label{exo:b1:structures:5}
设 $f \colon (\R, +) \to (\C^*, \times)$，$\theta \mapsto
\eu^{\iu\theta}$。证明 $f$ 是\omterm{def:b1:structures:morphism}{同态}，求 $\ker f$ 与
$\operatorname{im} f$，并由
\cref{prop:b1:structures:kernel} 推出 $f$ 不是\omterm{def:b1:logic:inj}{单射}。限制
定义域，使其在尽可能大的区间上为\omterm{def:b1:logic:inj}{单射}。
\end{exercise}

\begin{exercise}[$\star\star$]\label{exo:b1:structures:6}
设 $H, K$ 为 $G$ 的\omterm{def:b1:structures:subgroup}{子群}。证明 $H \cap K$ 是\omterm{def:b1:structures:subgroup}{子群}，
并证明 $H \cup K$ 是\omterm{def:b1:structures:subgroup}{子群}，\emph{仅当} $H \subseteq K$ 或
$K \subseteq H$。\emph{（若 $h \in H \setminus K$，$k \in K
\setminus H$，$hk$ 可能落在哪里？）}
\end{exercise}

\begin{exercise}[$\star\star$]\label{exo:b1:structures:7}
群 $G$ 满足 $x^2 = e$，对所有 $x \in G$ 成立。证明 $G$ 是
\omterm{def:b1:structures:group}{阿贝尔群}。\emph{（展开 $(xy)^2$。）}
\end{exercise}

\begin{exercise}[$\star\star$]\label{exo:b1:structures:8}
在 $\Z/18\Z$ 中：列出全部可逆元，并求 $\overline 5$ 的逆元；
解方程 $\overline 5\, x = \overline 7$；解方程 $\overline 6\, x =
\overline 3$ 与 $\overline 6\, x = \overline{12}$。
\end{exercise}

\begin{exercise}[$\star\star$]\label{exo:b1:structures:9}
证明\omterm{def:b1:logic:sets}{集合} $\Z[\sqrt 2] = \{a + b\sqrt 2 : a, b \in \Z\}$ 是
一个环（$\R$ 的子环），并证明 $1 + \sqrt 2$ 是其中的一个可逆元，
其各次幂两两不同，有无穷多个——因此 $\Z[\sqrt 2]^\times$
是无限集，与 $\Z^\times$ 不同。
\end{exercise}

\begin{exercise}[$\star\star\star$]\label{exo:b1:structures:10}
（布尔环）设 $A$ 为一个环，其中 $x^2 = x$ 对每个 $x$ 都成立。
证明 $x + x = 0$ 对所有 $x$ 成立，并证明 $A$ 是交换环。
\emph{（展开 $(x+x)^2$ 与 $(x+y)^2$。）} 给出这类环的一个例子：
取 $\mathcal{P}(E)$ 为底集，以对称差为加法、
以交为乘法。
\end{exercise}

\begin{exercise}[$\star\star\star$]\label{exo:b1:structures:11}
设 $G$ 为一个群，其中对某个固定的 $n \geq 1$，有 $(xy)^n =
x^n y^n$、$(xy)^{n+1} = x^{n+1}y^{n+1}$ 与 $(xy)^{n+2} =
x^{n+2}y^{n+2}$ 对所有 $x, y$ 成立。证明 $G$ 是\omterm{def:b1:structures:group}{阿贝尔群}。
\emph{（由三个恒等式出发，先推出 $y^n x = x y^n$，再
推出 $y^{n+1} x = x y^{n+1}$，并得出结论。）}
\end{exercise}

\begin{exercise}[$\star\star$]\label{exo:b1:structures:12}
\begin{enumerate}
  \item 求从 $(\Z, +)$ 到 $(\Z, +)$ 的所有群\omterm{def:b1:structures:morphism}{同态}。
  \item 证明从 $(\Q, +)$ 到
        $(\Z, +)$ 的唯一群\omterm{def:b1:structures:morphism}{同态}是零\omterm{def:b1:structures:morphism}{同态}。\emph{（对 $x \in \Q$ 及
        $n \in \N^*$，比较 $f(x)$ 与 $n\,f(x/n)$。）}
\end{enumerate}
\end{exercise}

\section{周末问题：对称群与八数码}

\begin{problem}\label{pb:b1:structures:1}
群 $\mathfrak S_n$（$\intint1n$ 上的\omterm{def:b1:counting:objects}{置换}群）是
数学中最古老的群，至今仍最富启发性。本题
从零开始建立其结构理论——包括轮换、
由对换生成，以及\emph{符号差}\omterm{def:b1:structures:morphism}{同态}
$\varepsilon \colon \mathfrak S_n \to \{\pm1\}$（其存在性
确有难度），以及\omterm{pb:b1:structures:1}{交错群} $\mathfrak A_n$
由 $3$-轮换生成——再将其用于一个经典谜题：
在 $3 \times 3$ 的滑块游戏中，
任何移动序列都不能只交换两块而保持其余不动。
\index{对称群}\index{符号差}\index{交错群}
\omterm{def:b1:counting:objects}{置换}作用于 $\intint1n$ 上；乘积 $\sigma\tau$ 表示
“先作用 $\tau$”；$[\,v_1, \dots, v_n]$ 表示
把 $i$ 映为 $v_i$ 的\omterm{def:b1:counting:objects}{置换}。

\medskip
\noindent\textbf{第一部分——轮换与对换。}
\begin{enumerate}
  \item 说明 $\abs{\mathfrak S_n} = n!$
        （由\cref{thm:b1:counting:counts}可知）。在 $\mathfrak S_3$ 中，
        计算 $\sigma = [2, 3, 1]$ 与 $\tau =
        [1, 3, 2]$ 的两种乘积，并得出 $\mathfrak S_3$ 不是
        \omterm{def:b1:structures:group}{阿贝尔群}。
  \item 一个 \emph{$k$-轮换} $(a_1\ a_2\ \dots\ a_k)$（$k \geq 2$，
        $a_i$ 两两不同）将 $a_1 \mapsto a_2
        \mapsto \dots \mapsto a_k \mapsto a_1$，并固定
        其余元素；其\emph{支集}为 $\{a_1, \dots,
        a_k\}$。证明支集不相交的两个轮换
        可交换。
  \item 证明每个 $\sigma \in \mathfrak S_n$ 都是
        支集两两不相交的轮换之积，并证明该分解
        在不计因子次序的意义下唯一。
        \emph{（对每个 $i$，考虑序列 $i, \sigma(i),
        \sigma^2(i), \dots$：它必回到 $i$；所得的
        \emph{轨道}\omterm{thm:b1:logic:partition}{划分} $\intint1n$，且 $\sigma$ 在
        每个轨道上都作为一个轮换。）}
  \item 将 $\sigma = [4, 1, 5, 2, 3, 7, 8, 6] \in
        \mathfrak S_8$ 分解为互不相交的轮换。
        依 $\sigma$ 的\emph{阶}（见
        \cref{def:b1:structures:order}），证明不相交轮换之积的阶
        等于各轮换长度的\omterm{def:b1:arith:lcm}{最小公倍数}，
        并求此 $\sigma$ 的阶。
  \item 证明逐项相消恒等式
        \[
        (a_1\ a_2\ \dots\ a_k)
        = (a_1\ a_k)(a_1\ a_{k-1})\cdots(a_1\ a_2) ,
        \]
        并由此得出每个\omterm{def:b1:counting:objects}{置换}都是
        对换之积。把第 4 小题中的 $\sigma$ 写成这种
        乘积。
  \item 进一步证明：\emph{相邻}对换
        $(i\ \ i{+}1)$ 已足够：对 $a < b$，
        \[
        (a\ b) = (a\ \ a{+}1)(a{+}1\ \ a{+}2)\cdots(b{-}1\ \ b)
        \cdots(a{+}1\ \ a{+}2)(a\ \ a{+}1),
        \]
        它是 $2(b - a) - 1$ 个相邻对换之积——
        一个\emph{奇数}（这一奇偶性在
        下文中将两次起作用）。
\end{enumerate}

\medskip
\noindent\textbf{第二部分——符号差的存在性。}对 $\sigma
\in \mathfrak S_n$，令
\[
N(\sigma) = \#\bigl\{(i, j) : i < j,\ \sigma(i) >
\sigma(j)\bigr\}
\]
为其\emph{逆序对}的个数，并令 $\varepsilon(\sigma) =
(-1)^{N(\sigma)}$。
\begin{enumerate}[resume]
  \item 计算 $N$ 与 $\varepsilon$ 在恒等\omterm{def:b1:counting:objects}{置换}上的值，
        对对换 $(i\ \ i{+}1)$ 及 $[2, 3, 1]$ 同样计算。
  \item 证明：对每个 $\sigma$ 及每个相邻对换
        $\tau = (i\ \ i{+}1)$，有
        $N(\sigma\tau) = N(\sigma) \pm 1$。\emph{（在右侧与
        $\tau$ 复合时，交换位置 $i$ 与
        $i + 1$ 上的值；恰有一对元素
        的逆序状态改变。）}
  \item 利用第 6 小题，推出对\emph{任意}
        对换 $\tau$，有 $\varepsilon(\sigma\tau) =
        -\varepsilon(\sigma)$；由此得出：若 $\sigma$ 是
        $p$ 个对换之积，则 $\varepsilon(\sigma)
        = (-1)^p$——特别地，$p$ 的奇偶性只依赖于
        $\sigma$，与所取的分解无关——并且
        $\varepsilon \colon \mathfrak S_n \to \{\pm 1\}$ 是一个
        群\omterm{def:b1:structures:morphism}{同态}。
  \item 证明 $k$-轮换的符号差为 $(-1)^{k-1}$，
        并证明一般地有 $\varepsilon(\sigma) = (-1)^{n -
        c(\sigma)}$，其中 $c(\sigma)$ 为
        $\sigma$ 的轨道个数（计入不动点）。
  \item \emph{\omterm{pb:b1:structures:1}{交错群}}为 $\mathfrak A_n =
        \ker\varepsilon$。说明它是\omterm{def:b1:structures:subgroup}{子群}，并证明
        $\abs{\mathfrak A_n} = \frac{n!}2$（$n \geq 2$）。
        \emph{（固定一个对换 $\tau_0$，考虑 $\sigma
        \mapsto \sigma\tau_0$。）}
  \item 对 $\sigma = [4, 1, 5, 2, 3, 7, 8, 6]$ 作一致性检验：
        用三种方法计算 $\varepsilon(\sigma)$——通过计数逆序对，
        利用第 10 小题由轮换类型计算，
        以及由第 5 小题中所得对换个数计算。
\end{enumerate}

\medskip
\noindent\textbf{第三部分——$\mathfrak A_n$ 由
$3$-轮换生成。}
\begin{enumerate}[resume]
  \item 设 $a, b, c, d$ 两两不同，验证下列两个
        恒等式
        \[
        (a\ b)(a\ c) = (a\ c\ b),
        \qquad
        (a\ b)(c\ d) = (a\ c\ b)(a\ c\ d) .
        \]
  \item 证明：当 $n \geq 3$ 时，$\mathfrak
        A_n$ 的每个元素都是 $3$-轮换之积。\emph{（偶\omterm{def:b1:counting:objects}{置换}
        可写成偶数个互换之积；
        每次两两一组地吸收。）}
  \item 将 $(1\ 2)(3\ 4)$ 与 $5$-轮换 $(1\ 2\ 3\ 4\ 5)$
        显式写成 $3$-轮换之积。
  \item 证明共轭公式：对每个 $\sigma \in
        \mathfrak S_n$,
        \[
        \sigma\,(a_1\ \dots\ a_k)\,\sigma^{-1}
        = \bigl(\sigma(a_1)\ \dots\ \sigma(a_k)\bigr) .
        \]
\end{enumerate}

\medskip
\noindent\textbf{第四部分——八数码问题。} 方块 $1, \dots, 8$ 在
$3 \times 3$ 的框中滑动，框内有一个空格；\emph{移动}是指
将与空格相邻的一个方块滑入空格。将格子依次编号为
$1, \dots, 9$（逐行；已解局面中，方块 $i$ 位于
格 $i$，空格位于格 $9$）。将空格视作
第九个方块，于是局面即为\omterm{def:b1:counting:objects}{置换} $\sigma \in \mathfrak
S_9$（方块 $\sigma(i)$ 位于格 $i$）。
\begin{enumerate}[resume]
  \item 证明一次移动将 $\sigma$ 换为 $\sigma \circ \tau$，
        其中 $\tau$ 为所涉两格的互换；
        由此推出每次移动都使
        $\varepsilon(\sigma)$ 变号。
  \item 设 $d(\sigma)$ 为曼哈顿距离（行距与
        列距之和），即空格当前所在格与其归位格
        $9$ 之间的距离。证明每次移动使 $d$ 改变 $\pm1$，
        从而每次移动也使 $(-1)^{d(\sigma)}$ 变号。由此可得
        \[
        I(\sigma) = \varepsilon(\sigma)\cdot(-1)^{d(\sigma)}
        \]
        在每次移动下都是\emph{不变量}。
  \item 证明该谜题的经典不可能性：
        将方块 $7$ 与 $8$ 互换、其余一切（包括空格）保持原位的局面，
        无法从已解局面
        到达。
  \item 在本阶段不加证明地承认其逆命题（其证明是一个有启发性但
        冗长的归纳法）：$I = +1$ 的每个局面均
        可到达。由此可得空格位于归位处的 $8!$ 个局面中
        恰有一半可解，即\
        $\frac{8!}2 = 20\,160$.
  \item 由第 20 问可知，空格位于归位处时，可达的方块排列
        恰好构成
        \omterm{def:b1:structures:subgroup}{子群} $\mathfrak A_8 \leq \mathfrak S_8$。
  \item 不变量的应用：能否到达 (a) 这样的局面：
        方块 $1, 2, 3$ 被循环\omterm{def:b1:counting:objects}{置换}，
        其余一切（包括空格）均归位？
        (b) 能否到达这样的局面：方块 $5$ 与空格
        互换了位置，其余方块均归位？
        请用 $I$ 分别说明两个答案的理由。
\end{enumerate}

\medskip
\noindent\textbf{第五部分——综合。}
\begin{enumerate}[resume]
  \item 证明：当 $n \geq 3$ 时，群\omterm{def:b1:structures:morphism}{同态} $f
        \colon \mathfrak S_n \to \{\pm 1\}$ 仅有两个：常值\omterm{def:b1:structures:morphism}{同态}
        与 $\varepsilon$。\emph{（利用第 16 问及
        $\{\pm1\}$ 的交换性，证明 $f$ 在所有互换上
        取值相同。）}
  \item 本题究竟在何处用到了：(i) \omterm{def:b1:structures:morphism}{同态}概念
        及\cref{prop:b1:structures:kernel}；(ii) 在
        \cref{ch:b1:counting}中的计数原理；(iii) 第 8--9 问所解决的
        良定义问题？每项各用
        一句话回答。
  \item 以一段短文作综述：一个奇偶性函数，
        良定义性只需验证一次，便能同时刻画
        $\mathfrak S_n$ 的内部结构（\omterm{def:b1:structures:subgroup}{子群}
        $\mathfrak A_n$），判定一个实物谜题，并且——
        借助公式 $\det = \sum_\sigma
        \varepsilon(\sigma)\cdots$——将
        在\cref{ch:b1:det}中定义行列式。评述这一反复出现的模式：
        不变量把“穷举一切移动序列”化为一次
        计算。
\end{enumerate}
\end{problem}
